% 【本文件写什么】impl 实现层：dummy
%   - 定位：占位/链接桩实现，保证默认 feature 可编译。
%   - crate 路径：os/components/wateros-mm/mm-frame-alloctor/frame-alloctor-impl/impl-dummy/src/
%   - 如何实现对应 api-v0 trait：关键类型、状态机、数据结构
%   - 算法或硬件交互流程（可用伪代码/流程图）
%   - 本 impl 启用的 Cargo [features] 与 arch feature（impl-riscv64 等）
%   - 与其它 impl 变体的差异、切换 feature 时的影响
%   - 性能/内存假设与已知限制
% 【事实来源】
%   - 上述 crate 源码与 Cargo.toml
%   - os/feature-tree.txt 中指向本 impl 的 feature 链

\subsubsection{impl-dummy（帧分配）}

\code{wateros-mm-frame-alloctor-impl-dummy}为空 crate：不导出\code{init\_frame\_allocator}、\code{frame\_alloc\_result}等全局符号。供仅编译\code{mm-api}契约、或组合\code{impl-dummy}页表桩的 cfg 使用。

任何依赖真实物理帧的路径（页表映射、COW、惰性缺页）须在 Cargo feature 中启用\code{impl-stack}。与\code{impl-stack}切换时，\code{frame-alloctor}聚合\code{lib.rs}通过\code{cfg(feature = "impl-stack")}选择再导出目标。
